\section{Symptom Token representation}
\label{sec:tokens}

\begin{definition}[Symptom Token]
For each aggregated channel $j$ in $\{1,\ldots,22\}$ (labs + vitals), a \textbf{Symptom Token} is a tuple
\[
  \phi_j = \bigl(\texttt{symptom\_id}_j,\; \texttt{intensity}_j \in [0,1],\; \texttt{time\_delta}_j,\; \texttt{modality}_j \in \{\text{lab},\text{vital}\}\bigr),
\]
where \texttt{intensity} is min--max normalized after clipping to clinically plausible ranges; \texttt{time\_delta} encodes position within the 24\,h window; missing channels are handled per repository imputation rules.
\end{definition}

\subsection{Alignment with tabular branch}
The same raw measurements produce $\mathbf{x} \in \mathbb{R}^{25}$ (with engineered summaries). Thus neural and classical models consume \textbf{mutually consistent} physiology, enabling fair fusion and ablation (see Exp4 in companion tables: removing symptom tokens changes the pathway).

\subsection{Why tokens?}
Self-attention~\citep{vaswani2017attention} can learn pairwise interactions among channels without hand-specified crosses; tokens also interface with few-shot episodic training that expects sequence-like structure.

\begin{figure}[htbp]
  \centering
  \resizebox{\textwidth}{!}{% MIMIC-IV extraction schema — monochrome, publication-style (TikZ)
% Preamble (once per document): \usepackage{tikz}\usetikzlibrary{positioning,arrows.meta,calc}
\begin{tikzpicture}[
  scale=0.88,
  transform shape,
  font=\footnotesize,
  >=Stealth,
  tbl/.style={
    rectangle,
    draw=black!55,
    line width=0.4pt,
    fill=black!3,
    minimum height=8mm,
    inner xsep=4pt,
    inner ysep=3pt,
    align=center,
    rounded corners=1.5pt,
  },
  code/.style={tbl, font=\footnotesize\ttfamily},
  proc/.style={
    rectangle,
    draw=black!55,
    line width=0.45pt,
    fill=black!6,
    inner sep=6pt,
    align=center,
    rounded corners=2pt,
    font=\footnotesize,
  },
  outbox/.style={
    rectangle,
    draw=black!55,
    line width=0.45pt,
    fill=black!10,
    inner sep=6pt,
    align=center,
    rounded corners=2pt,
    font=\footnotesize\bfseries,
  },
  ed/.style={black!50, line width=0.35pt},
  el/.style={font=\footnotesize, text=black!65, inner sep=1pt, fill=white},
]

% --- Core hospital tables ---
\node[code] (pat) {patients};
\node[code, right=9mm of pat] (adm) {admissions};
\node[code, right=9mm of adm] (icu) {icustays};
\node[code, above=9mm of icu, xshift=-6mm] (dx) {diagnoses\_icd};

\draw[ed, ->] (pat) -- node[el, above=0.5pt] {\texttt{subject\_id}} (adm);
\draw[ed, ->] (adm) -- node[el, above=0.5pt] {\texttt{hadm\_id}} (icu);
\draw[ed, ->] (adm.north) |- node[el, pos=0.28, left=1pt] {\texttt{hadm\_id}} (dx.west);

% --- Labs / chart (second row) ---
\node[code, below left=12mm and 2mm of icu] (dlab) {d\_labitems};
\node[code, right=5mm of dlab] (lab) {labevents};
\node[code, right=7mm of lab] (chart) {chartevents};
\node[code, right=5mm of chart] (dit) {d\_items};

\draw[ed, ->] (dlab) -- node[el, above=0.5pt] {\texttt{itemid}} (lab);
\draw[ed, ->] (dit) -- node[el, above=0.5pt] {\texttt{itemid}} (chart);

\coordinate (hub) at ($(icu.south)+(0,-2.5mm)$);
\draw[ed, ->] (icu.south) -- (hub);
\draw[ed, ->] (hub) -| node[el, pos=0.18, below=1pt] {\texttt{stay\_id}} (lab.north);
\draw[ed, ->] (hub) -| (chart.north);

% --- Preprocessing ---
\node[proc, below=11mm of lab, anchor=north, xshift=16mm] (prep) {%
  \textbf{Preprocessing pipeline}\\[0.35em]
  {\footnotesize 15 labs + 7 vitals $\rightarrow$ clip / min--max $[0,1]$ $\rightarrow$ feature engineering}\\[0.2em]
  {\footnotesize $\rightarrow$ \textbf{25-D tabular} $\;|\;$ \textbf{token sequence} $\rightarrow$ ETHOS $\rightarrow$ \textbf{128-D}}%
};

\coordinate (merge) at ($(prep.north)+(0,4.5mm)$);
\draw[ed, ->] (lab.south) |- (merge);
\draw[ed, ->] (chart.south) |- (merge);
\draw[ed, ->] (merge) -- (prep.north);

% --- Cohort output ---
\node[outbox, below=6mm of prep] (cohort) {$2{,}000$ patients $\times$ $25$ features + labels};
\draw[ed, ->] (prep) -- (cohort);

\end{tikzpicture}
}
  \caption{From MIMIC tables to tabular features and token sequences (vector schema, monochrome).}
  \label{fig:schema_c}
\end{figure}
